% ---
% title: 富比尼定理
% date: 2026/2/26
% tag: Math
% ---
\textbf{富比尼定理}（英語：Fubini's theorem）是\href{https://zh.wikipedia.org/wiki/数学分析}{數學分析}中有關重積分的一個定理，由數學家\href{https://zh.wikipedia.org/wiki/圭多·富比尼}{圭多·富比尼}在1907年提出。富比尼定理給出了使用\href{https://zh.wikipedia.org/wiki/逐次积分}{逐次積分}的方法計算\href{https://zh.wikipedia.org/wiki/双重积分}{雙重積分}的條件。在這些條件下，不僅能夠用逐次積分計算雙重積分，而且交換逐次積分的順序時，積分結果不變。「托內利定理」由數學家\href{https://zh.wikipedia.org/wiki/列奧尼達·托內利}{列奧尼達·托內利}在1909年提出，與富比尼定理相似，但是是應用於非負函數而不是可積函數。\href{https://zh.wikipedia.org/wiki/富比尼定理#cite_note-1}{[1]}

\section{定理}

若

\[
{\displaystyle \int _{A\times B}|f(x,y)|\,d(x,y)<\infty }
\]

其中\emph{A}和\emph{B}都是\href{https://zh.wikipedia.org/wiki/Σ-有限测度}{σ-有限測度}空間，\( {\displaystyle A\times B}\)是\(A\)和\(B\)的\href{https://zh.wikipedia.org/wiki/积测度}{積可測空間}，\({\displaystyle f:A\times B\mapsto \mathbb {C} }\)是\href{https://zh.wikipedia.org/wiki/可测函数}{可測函數}，那麼

\[
{\displaystyle \int _{A}\left(\int _{B}f(x,y)\,dy\right)\,dx=\int _{B}\left(\int _{A}f(x,y)\,dx\right)\,dy=\int _{A\times B}f(x,y)\,d(x,y)}
\]

前二者是在兩個測度空間上的\href{https://zh.wikipedia.org/wiki/逐次积分}{逐次積分}，但積分次序不同；第三個是在乘積空間上關於乘積測度的積分。

特別地，如果\({\displaystyle f(x,y)=h(x)g(y)}\)，則

\[
{\displaystyle \int _{A}h(x)\,dx\int _{B}g(y)\,dy=\int _{A\times B}f(x,y)\,d(x,y)}
\]

如果條件中絕對值積分值不是有限，那麼上述兩個逐次積分的值可能不同。

\section{非負函數的Tonelli定理}

\href{https://zh.wikipedia.org/wiki/列奧尼達·托內利}{托內利}定理延續了富比尼定理。托內利定理的結論與富比尼定理一樣，但是條件從 \({\displaystyle |f|}\) 積分有限改為了\({\displaystyle f}\) 非負。

\section{應用}

富比尼定理一個應用是計算\href{https://zh.wikipedia.org/wiki/高斯积分}{高斯積分}。高斯積分是很多概率論結果的基礎：

\[
{\displaystyle \int _{-\infty }^{\infty }e^{-\alpha x^{2}}\,dx={\sqrt {\frac {\pi }{\alpha }}}}
\]

\section{參考文獻}

\begin{enumerate}
  \item \href{https://zh.wikipedia.org/wiki/列奧尼達·托內利}{Tonelli, Leonida} (1909). "Sull'integrazione per parti". \emph{Atti della Accademia Nazionale dei Lincei}. (5). \textbf{18} (2): 246–253.
\end{enumerate}
